Los dos casos frontera más citados en este informe -- la comparación
0{,}5\,ms vs.\ 1\,ms de \texttt{interval\_ns}, y la comparación 50\,ms vs.\
100\,ms de \texttt{gpu\_interval\_ns} en \texttt{rodinia\_backprop} -- se
re-midieron con las 24 combinaciones intercaladas en \textbf{orden
aleatorio} (semilla fija, reproducible), en vez del orden fijo por bloques
de los barridos principales, para descartar que una deriva temporal del
nodo compartido explicara el resultado en vez del parámetro mismo.

\begin{table}[H]
\centering
\caption{CPU: CV\% del intervalo real de muestreo, orden aleatorizado (3 repeticiones c/u).}
\label{tab:randomized-cpu}
\small
\begin{tabular}{@{}lrr@{}}
\toprule
\textbf{Kernel} & \textbf{0{,}5\,ms (media)} & \textbf{1\,ms (media)} \\
\midrule
\texttt{npb\_bt} & 0{,}096\,\% & 0{,}061\,\% \\
\texttt{npb\_cg} & 0{,}146\,\% & 0{,}095\,\% \\
\texttt{npb\_mg} & 0{,}284\,\% & 0{,}191\,\% \\
\bottomrule
\end{tabular}
\end{table}

Los tres kernels reproducen, bajo orden aleatorizado, la misma dirección
observada en el barrido principal (0{,}5\,ms sistemáticamente más ruidoso
que 1\,ms) -- no hay indicio de que el orden fijo original haya introducido
un sesgo que infle artificialmente esta diferencia.

\begin{table}[H]
\centering
\caption{GPU: muestras \texttt{gpu\_telemetry}, \texttt{rodinia\_backprop}, orden aleatorizado.}
\label{tab:randomized-gpu}
\small
\begin{tabular}{@{}lrr@{}}
\toprule
\textbf{Repetición (orden real de ejecución)} & \textbf{Cadencia} & \textbf{Muestras} \\
\midrule
1 & 50\,ms  & 17 \\
2 & 50\,ms  & 17 \\
3 & 100\,ms & 9 \\
4 & 100\,ms & 9 \\
5 & 100\,ms & 9 \\
6 & 50\,ms  & 17 \\
\bottomrule
\end{tabular}
\end{table}

El caso GPU es aún más contundente: bajo intercalado aleatorio, las tres
corridas de 50\,ms dan exactamente 17 muestras y las tres de 100\,ms dan
exactamente 9 -- una reproducción exacta, cifra por cifra, de los valores
del barrido principal no aleatorizado (Sección~\ref{sec:gpu-interval-ns}).

\textbf{Conclusión}: para estos dos casos frontera, la ausencia de
aleatorización en el diseño original de los barridos no parece haber
introducido un sesgo material -- el efecto atribuido al parámetro se
reproduce bajo un orden de ejecución distinto. Esto \textbf{no} es una
prueba general de que ningún barrido de este informe esté libre de deriva
temporal del nodo compartido (los otros 20 kernels y las otras cadencias no
se reconfirmaron de esta forma, por costo de cómputo), pero sí reduce el
riesgo señalado para los dos casos donde más importaba, y establece que el
método de confirmación dirigida es una vía barata y disponible si en el
futuro se quisiera extenderlo a más casos.